\documentclass[11pt]{article}

\usepackage[margin=1in]{geometry}
\usepackage[T1]{fontenc}
\usepackage{lmodern}

\usepackage{amsmath, amssymb}

\usepackage{setspace}

\usepackage[most]{tcolorbox}
\usepackage{xcolor}

\tcbset{
  enhanced,
  boxrule=0pt,
  frame hidden,
  breakable
}

\usepackage{fancyhdr}
\pagestyle{fancy}
\fancyhf{}
\setlength{\headheight}{52.42322pt}
\renewcommand{\headrulewidth}{0pt}

\newcommand{\Course}{Math 4610 Spring 2026}
\newcommand{\Instructor}{Discussion 1}
\newcommand{\AssignmentType}{Discussion} % or Discussion
\newcommand{\AssignmentNumber}{1}

\newcommand{\Contributors}{
Marks, Stephen\\
}

\fancyhead[L]{\Course\\ \Instructor}
\fancyhead[R]{\Contributors}

\newcommand{\ProblemDivider}[1]{%
  \vspace{1.0em}
  \noindent\textbf{#1}\par
  \vspace{0.35em}
}

\definecolor{problemBack}{RGB}{235,242,250}
\definecolor{problemFrame}{RGB}{70,110,160}
\definecolor{exerciseGray}{gray}{0.96}
\definecolor{exerciseGrayFrame}{gray}{0.45}

\newtcolorbox{problemBox}{
  colback=exerciseGray,
  colframe=exerciseGrayFrame,
  arc=2mm,
  left=8pt,right=8pt,top=8pt,bottom=8pt
}

\newtcolorbox{solutionBox}{
  colback=white,
  colframe=white,
  arc=0mm,
  left=0pt,right=0pt,top=0pt,bottom=0pt
}

\newcommand{\ProblemAnnounce}{\noindent}
\newcommand{\SolutionAnnounce}{\noindent\textbf{Solution.}\par\medskip}

\newcommand{\Exercise}[1]{\textbf{Exercise~#1.}\;}

\newcommand{\R}{\mathbb{R}}
\newcommand{\C}{\mathbb{C}}

\begin{document}

\begin{center}
  {\Large \textbf{\AssignmentType~\AssignmentNumber}}
\end{center}
\vspace{0.75em}

\ProblemDivider{}

\begin{problemBox}
\ProblemAnnounce
\Exercise{1}
Let $m,n$ be two positive integers. Define
\[
\mathbb{F}^{m,n}:=\{\text{$m\times n$ matrices with entries in $\mathbb{F}$}\}.
\]
See Definition 3.39 in Section 3C. Suppose $T\in \mathcal{L}(V)$, and $\{u_1,\dots,u_n\}$ and $\{v_1,\dots,v_n\}$ are bases of $V$.
Prove that the following are equivalent.
\begin{enumerate}
  \item[(a)] $T$ is invertible.
  \item[(b)] The columns of $\mathcal{M}(T)$ are linearly independent in $\mathbb{F}^{n,1}$.
  \item[(c)] The columns of $\mathcal{M}(T)$ span $\mathbb{F}^{n,1}$.
  \item[(d)] The rows of $\mathcal{M}(T)$ are linearly independent in $\mathbb{F}^{1,n}$.
  \item[(e)] The rows of $\mathcal{M}(T)$ span $\mathbb{F}^{1,n}$.
\end{enumerate}
Here $\mathcal{M}(T)$ means $\mathcal{M}\!\bigl(T,(u_1,\dots,u_n),(v_1,\dots,v_n)\bigr)$ (see Definition 3.31 in Section 3C).
You can refer to \emph{Rank of Matrix} (page 79) in Section 3C. (Do not use determinants.)
\end{problemBox}

\begin{solutionBox}
\SolutionAnnounce
\doublespacing

\subsection*{(a $\to$ b)}
Assume (a) and suppose that $\mathcal{M}(T)$ has entries $A_{j,k}$ 
Let $b_1,\dots,b_n\in\mathbb{F}$ be scalars such that
\[
b_1\begin{pmatrix}A_{1,1}\\ \vdots\\ A_{n,1}\end{pmatrix}
+\cdots+
b_n\begin{pmatrix}A_{1,n}\\ \vdots\\ A_{n,n}\end{pmatrix}
=
\begin{pmatrix}b_1A_{1,1}+\cdots+b_nA_{1,n}\\ \vdots\\ b_1A_{n,1}+\cdots+b_nA_{n,n}\end{pmatrix}
=
\begin{pmatrix}0\\ \vdots\\ 0\end{pmatrix}.
\]
Let $u=b_1u_1+\cdots+b_nu_n$ so that
\begin{align*}
Tu
&= b_1Tu_1+\cdots+b_nTu_n \\
&= b_1(A_{1,1}v_1+\cdots+A_{n,1}v_n)+\cdots+b_n(A_{1,n}v_1+\cdots+A_{n,n}v_n)\\
&= (b_1A_{1,1}+\cdots+b_nA_{1,n})v_1+\cdots+(b_1A_{n,1}+\cdots+b_nA_{n,n})v_n\\
&=0v_1+\cdots+0v_n=0.
\end{align*}
Thus $u\in \operatorname{null}T$. Since $T$ is invertible, and thus injective, $u=0$.
The linear independence of the basis $u_1,\dots,u_n$ then gives $b_1=\cdots=b_n=0$.
Thus the columns of $\mathcal{M}(T)$ are linearly independent.

\subsection*{(b $\to$ a)}
Now assume (b) and let $u=b_1u_1+\cdots+b_nu_n$ be such that $Tu=0$.
Like above, this is equivalent to
\[
(b_1A_{1,1}+\cdots+b_nA_{1,n})v_1+\cdots+(b_1A_{n,1}+\cdots+b_nA_{n,n})v_n=0.
\]
Due to linear independence of the basis $v_1,\dots,v_n$,
\[
b_1A_{1,1}+\cdots+b_nA_{1,n}=\cdots=b_1A_{n,1}+\cdots+b_nA_{n,n}=0,
\]

\[
\begin{pmatrix}b_1A_{1,1}+\cdots+b_nA_{1,n}\\ \vdots\\ b_1A_{n,1}+\cdots+b_nA_{n,n}\end{pmatrix}
=
b_1\begin{pmatrix}A_{1,1}\\ \vdots\\ A_{n,1}\end{pmatrix}
+\cdots+
b_n\begin{pmatrix}A_{1,n}\\ \vdots\\ A_{n,n}\end{pmatrix}
=
\begin{pmatrix}0\\ \vdots\\ 0\end{pmatrix}.
\]
$b_1=\cdots=b_n=0$, because the columns of $\mathcal{M}(T)$ are linearly independent. So $u=0$.
Thus $T$ is injective. Because $T\in\mathcal{L}(V)$ with $\dim V=n<\infty$, injectivity implies invertibility, so $T$ is invertible.

\subsubsection*{Therefore, (a $\leftrightarrow$ b)}

\subsection*{(b $\to$ c)}
Assume (b). Since the columns are linearly independent in $\mathbb{F}^{n,1}$, and are a list of $n$ vectors, they span $\mathbb{F}^{n,1}$ since $\dim(\mathbb{F}^{n,1})=n$.

\subsection*{(c $\to$ b)}
Assume (c). Since the columns span $\mathbb{F}^{n,1}$ and are a list of $n$ vectors in an $n$-dimensional space, they form a basis. So they are linearly independent in $\mathbb{F}^{n,1}$ by the definition of a basis.

\subsubsection*{Therefore, (b $\leftrightarrow$ c)}

\subsection*{(c $\to$ e)}
Recall:
\begin{quote}
3.57: Suppose $A\in \mathbb{F}^{m,n}$. Then the column rank of $A$ equals the row rank of $A$.
\end{quote}
Assume (c). Then the column space of $\mathcal{M}(T)$ is equal to $F^{n,1}$, and thus
\[
  \dim(\operatorname{Col}(\mathcal{M}(T))) = n,
\]
so the column rank of  $\mathcal{M}(T)$ is $n$. By 3.57, the row rank of $\mathcal{M}(T)$ is also $n$. Since the row space of $\mathcal{M}(T)$ is a subspace of $F^{1,n}$, which has dimension $n$, it follows that the row space must be all of $F^{1,n}$. Thus, the rows of $\mathcal{M}(T)$ span $F^{1,n}$.

\subsection*{(e $\to$ c)}
Conversely, assume (e). Then the row space of $\mathcal{M}(T)$ is equal to $F^{1,n}$, so the row rank of $\mathcal{M}(T)$ is $n$. By 3.57, the column rank of $\mathcal{M}(T)$ is also $n$. Since the column space of $\mathcal{M}(T)$ is a subspace of $F^{n,1}$ with dimension $n$, it follows that the column space must be all of $F^{n,1}$. Thus, the columns of $\mathcal{M}(T)$ span $F^{n,1}$.
\subsubsection*{Therefore, (c $\leftrightarrow$ e)}

\subsection*{(d $\to$ e)}
Assume (d). Since the rows are linearly independent in $\mathbb{F}^{1,n}$ and are a list of $n$ vectors, they span $\mathbb{F}^{1,n}$ since $\dim(\mathbb{F}^{1,n})=n$.

\subsection*{(e $\to$ d)}
Assume (e). Since the rows span $\mathbb{F}^{1,n}$ and are a list of $n$ vectors in an $n$-dimensional space, they form a basis. So they are linearly independent in $\mathbb{F}^{1,n}$ by the definition of a basis.

\subsubsection*{Therefore, (d $\leftrightarrow$ e)}

So, by (a) $\leftrightarrow$ (b) $\leftrightarrow$ (c) $\leftrightarrow$ (e) $\leftrightarrow$ (d), (a) -- (e) are equivalent.


\end{solutionBox}

\hfill \ensuremath{\square}
\end{document}
